\documentclass[11pt]{article}
\usepackage{amsmath,amssymb,amsthm,geometry}\geometry{margin=1in}
\newtheorem{theorem}{Theorem}\newtheorem{proposition}[theorem]{Proposition}
\title{Whole-case terminal grades and exact event-free source recovery, v3}
\date{16 November 2026}
\begin{document}\maketitle
\paragraph{Audit and repair target.} Whole-case terminal-grade instantiation with an integrated event-free constructor/fold compatibility section. The prior v1 SHA256 \texttt{0e8f712886af71cbd4f06e4b3090f7e266468fcedb276ea48b720115588956c2} remains unchanged. The reviewed world-inert initializer-family theorem, SHA256 \texttt{b90dca1995c2a151ee75451ad1780dd70dee298b873b02578f99ab47fae79781}, has an operational proof independent of its terminal-grade recursion. Its original-source recovery paragraph requires the original \emph{whole uncolored case grade}, not merely the earlier genuinely slot-colored grade convention. This note supplies that separate functor/body/initial-algebra proof and the exact event-free recovery. The recovery supplied here requires this whole-case instantiation; the earlier operational theorem is unaffected.

\paragraph{Sources read.} Original ordered-polynomial source: \texttt{182a8462b3a7732967ad59f774dc2dd91090bf60c09a3329c9dd8c31e5bbd07b}. Original initial-row feedback/source recovery: \texttt{7ffa6f93b05b978f5529f8bf6f2ea79a5d3d34fb4afb1f57dbfa7c44e791da64}. The latter explicitly uses $C\circ(I+U\cdot C_X)$, not a realized-case restriction. The all-small full-family and finite actual coefficient assumptions remain in force.

\section{An independently specified whole-case grade functor}
Fix a primitive $P:A+U\rightsquigarrow B+U$ with entire original constant ordered-polynomial grade $C=(C_n,R_c)$. A completed root state supplies a shape $c\in C_n$ and $n$ ordered output values in $B+U$. The grade $C$ itself does not select, for every unused shape slot, whether its future value belongs to $B$ or $U$.

For an arbitrary closed target system $X:U\rightsquigarrow B$ with full grade $D$, define
\[
 G_{C,U}(D)=C\circ(I+U\cdot D),                         \tag{1}
\]
where $I$ has one unary shape with unit coefficient, and $U\cdot D$ has a copy of every shape of $D$ at each loop label $u$. This is the original full coproduct/case collection. A shape of (1) chooses $c\in C_n$ and, independently at each slot, either the $I$ alternative or $(u,d)$ with $u\in U$ and $d\in D$. Its external arity is the sum of those chosen alternatives' arities. Its actual coefficient is $R_c$ tensored with all chosen $D$ coefficients, retaining unit factors through the canonical unit maps. All choices occur, including choices incompatible with any actual root state.

An arity-preserving whole grade map $\gamma:D\to D'$ acts at every child $d$ by $\gamma$ and its actual coefficient arrow, leaving $c$, case tags and loop labels unchanged. Ordered substitution gives a genuine endofunctor with its usual actual coefficient coherence. No map $C\to C\circ(I+U\cdot D)$ is postulated on the basis of a state-dependent case choice.

The operational root constructor of the reviewed theorem is unchanged. On a completed state $(p,(x_j))$, inspect the actual body of $p$: at a $B$-valued slot use $I$ and retain that $B$ value; at a $U$-valued slot labeled $u$ use $(u,d_j)$ where $d_j$ is the full terminal body shape of the matching completed child, and substitute its ordered $B$ outputs. This gives a body in (1). It is supplied on every completed configuration, including unreachable ones. A native child map preserves the actual child full body and ordered ports, so this body commutes with $G_{C,U}(\gamma)$. Root/passive port stability and completed-body rebase invariance preserve the same case choice at actual states. Thus this is an actual full terminal-grade/body component of $T_P$, not only a carrier cardinality formula.

\section{The entire finite-tree initial algebra}
Let $R_UC$ have every nonempty finite rooted ordered $C$ tree in which each slot is either an external leaf or an internal edge labeled by $u\in U$ leading to a child $C$ tree. The shape's external positions are its ordered leaves. Its coefficient is the ordered tensor of all $R_c$ vertex factors; units at external slots are handled by actual unitors. This includes all formal trees, not only shapes realized by any execution or initializer law.

\begin{theorem}[Full whole-case initial algebra]
Root grafting $G_{C,U}(R_UC)\to R_UC$ is an isomorphism of entire weighted collections, and is an initial $G_{C,U}$ algebra in the arity-preserving weighted-collection category.
\end{theorem}
\begin{proof}
Root splitting uniquely recovers the root $c$ and the $I$ choice at each external slot or the loop label and child tree at each internal slot. It preserves the ordered external positions. Actual coefficients on both sides have exactly the same labeled vertex factors; the specified tensor associator/unit maps give inverse arrows. This proves the full isomorphism, including unused shapes.

For any actual grade algebra $a:G_{C,U}(D)\to D$, interpret each finite formal tree by recursively interpreting its child trees and applying $a$ to the resulting root case shape. Its coefficient arrow first tensors identity on $R_c$ with the recursively obtained child arrows, then composes with the actual coefficient arrow supplied by $a$. This defines arrows on every shape, with the correct ordered arity. The algebra equation holds by construction; finite-tree induction forces both shape and actual coefficient maps for any competing algebra morphism. No arrow into a terminal coefficient object is invented; it is supplied by $a$.\end{proof}

\begin{theorem}[Whole-case version of native feedback initiality]
Replace the terminal-grade component of the reviewed initializer-family category and constructor by (1) and its actual body above. On its same all-world all-raw well-founded initializer domain, the same finite operational feedback carrier with full terminal collection $R_UC$ is an initial native $T_P$ algebra in this whole-case category. Its command-indexed/global/flat writer branches retain exactly their independently stated hypotheses.
\end{theorem}
\begin{proof}
The operational raw/state/initial-family/expanded-paid-row and actual context constructions are unchanged. Their proofs do not identify scope writers with the terminal grade. For a target native algebra $\alpha$, its terminal component is now specifically an actual algebra $G_{C,U}(D)\to D$. The preceding theorem supplies the full grade fold on every formal tree and every coefficient. On completed source configurations, the actual root's $B/U$ output cases are exactly those used by the operational root split, so the grade-body equation follows by induction from the child body equations and the terminal body square of $\alpha$. This checks actual bodies rather than assuming that an arbitrary grade fold matches them. The independent operational fold, initial raw/coefficient recursion and paid successor equations remain those of the reviewed theorem. Uniqueness is their existing full-field induction together with the independently proved whole-case grade induction. Thus the proof is instantiated with a different specified grade functor, not with a silent equality between colored and uncolored grades.\end{proof}

\paragraph{Semantic target check.} The earlier no-external-output population algebra with only one nullary terminal shape is generally \emph{not} an algebra for (1): a formal $I$ alternative can have positive arity. In the unit-coefficient model, repair its terminal collection by taking one shape $d_n$ of unit coefficient for every finite arity $n$. Send every formal $G_{C,U}(D)$ shape of arity $n$ to $d_n$, using the actual unit-coefficient arrow. The population's completed states have zero external outputs and body $d_0$; the actual root body therefore maps to $d_0$ as required. All unused positive-arity shapes are mapped to their $d_n$ rather than discarded. The same population state/raw/schema algebra is then a meaningful opaque whole-case algebra. With arbitrary content coefficients, the required coefficient arrows must instead be supplied; no automatic augmentation in the independent content category is assumed.

\section{Exact original source recovery at the event-free interface}
Let $F\in H(A+U,B+U)$ be an original finite normalized ordered-polynomial process with whole grade $C$. Its zero-paid initial-row realization $J(F)$ has one completed state for each original raw witness, initializer selected-state map identity, original normalized initial law and actual augmented coefficient, and no paid commands. Assume its all-raw loop-output relation is well founded on all $U$, including zero-weight witnesses. This is the original exact finite silent domain, not almost-sure termination.

Every finite compatible invocation configuration of $J(F)$ is then exactly a raw initializer tree: each vertex state is its original raw witness and its emitted loop ports prescribe its children. Thus the family feedback object is itself event-free, its state and initial raw carriers coincide, and its selected-state map is identity. Its admitted law family is the recursively defined family of whole initializer-tree laws; its distinguished member is the original product of primitive probabilities. Its context coefficient at a tree is exactly the original tensor of all node context coefficients. Its full grade is now the entire $R_UC$ just proved, and its body is the original recursive tree rendering.

Let $D_H$ here simply mean selection of the supplied distinguished law, with no change to any raw, coefficient, whole-grade or body field. Then there is an actual native isomorphism
\[
 D_H\bigl(\mathrm{Fb}^{\mathsf F,case}_U(J(F))\bigr)
 \cong \operatorname{Rec}^{\mathrm{finite},\otimes}_U(F). \tag{2}
\]
The right side is the original finite well-founded invocation-tree recurrence. On loop inputs the isomorphism is identity on the complete raw trees and all actual factors; on external $A$ inputs use the same outer root and its loop-child trees. Full grade is identity on $R_UC$, not merely on realized terminal observations. Normalization and finite raw cardinality follow from the all-raw guard. There is no paid boot and no lost initial coefficient.

The comparison is natural for actual original native cells with fixed loop interface $U$ when both source and target satisfy the all-raw guard. Their raw initial maps, port-label bijections and actual coefficients map the same complete initializer trees; the entire original grade map acts at every vertex and hence on all of $R_UC$. Noninjective initial raw transport uses the positive-fiber convex-mixture argument for the family extension, while the distinguished product satisfies the original exact fiber equation. A map does not automatically transport the guard to a target with additional zero raw cycles; admission of both objects is a separate domain hypothesis.

The entire family object before applying $D_H$ may admit nonproduct correlations absent from the original singleton-law source, so (2) is not a full-family isomorphism to that source. For systems with paid menus, merely selecting the initial distinguished member does not restrict nonproduct newborn initializer commands, so this event-free theorem is not an unrestricted paid-runtime retraction. The independent scope-writer/terminal distinction also remains unchanged.

\paragraph{Colored versus whole-case branches.} Genuinely colored formal slots can yield a smaller grade and stronger post-tightening isomorphisms. They are not automatically identified with the whole-case collection above. A comparison requires a supplied full shape/coefficient map recording the chosen slot tags and loop labels on every unused shape; a state-dependent actual output value cannot supply such a global map. This note makes no automatic colored-to-whole-case equivalence and does not delete off-case shapes from (1).

\section{Integrated event-free constructor, fold and counit compatibility}
This section concerns the actual family mixed hom category $H_{\mathsf F}$, not a forgetful operation on paid systems. This v3 adds the following explicit correspondence; v2 SHA256 \texttt{7726da321d398a24495642fb2b74d4d481f5c879b6b49f51d54bf923118a0262} remains unchanged.

\paragraph{Exact event-free operational correspondence.} Work over the terminal world category, not a one-object category with nontrivial update endomorphisms. Take the command-indexed operational subcategory with no paid commands, all states completed, state carrier at each input equal to its finite common initial raw carrier, and initializer selected-state map the identity. Use canonical ordered rank ports and the constant whole-case terminal grade/body. There is no extra unused global paid-writer schema field in this subcategory. Its initializer family and distinguished member are exactly the $H_{\mathsf F}$ law data, and its actual initial coefficients are exactly the context coefficients. Passive rebase is the forced identity. Conversely these data define such an operational object from every $H_{\mathsf F}$ object. The initial selected-state square forces every operational state map to equal its common raw map; the ordered full body equation forces its rank-port bijection. Remaining fields are exactly the actual context arrows and whole grade/body map. These constructions are inverse on objects and native morphisms with canonical rank ports. If arbitrary stable port names are retained, the unique order-preserving name-to-rank bijections instead give a specified equivalence, not literal equality of presentations.

For event-free $P=J(F)$, the root constructor preserves this subcategory up to the specified forced-unit/port normalization: root states and child states are precisely the same tuples as root and child initial raw witnesses, so the decoder remains identity; all states stay completed and no paid commands appear. Pure $B$-identity copies contribute only forced singleton witnesses and unit factors. Its whole terminal collection and body are exactly (1), not the smaller colored collection. The all-raw initializer guard is the same relation on original raw output labels. Therefore the operational initiality theorem restricts to an actual $H_{\mathsf F}$ initial-algebra theorem through this fully specified correspondence; it does not require adding arbitrary opaque configurations to an $H_{\mathsf F}$ object. It uses the distinguished-law functor $D_H$ and singleton embedding $J_H$ with the exact fields described above. Their homwise adjunction is given in \emph{Distinguished Law Coreflection and Relative Recovery}, SHA256 \texttt{67c1f6b9498faeebb69dd5263f7eb4364b1cfdf210abecdf6b3c0bb5801fc470}: an arrow out of a singleton-law source has only its distinguished law to test, so the native Hom-set bijection retains all raw/coefficient/grade fields. The following assertions can also be checked directly from these fields.

Fix the original process $P:U\rightsquigarrow B+U$, with its normalized initial row and full grade $C$. For a family mixed process $X:U\rightsquigarrow B$, let $\operatorname{Case}_U^{\mathsf F}(X)$ use the singleton identity at a $B$ input and the corresponding $X$ row at a $U$ input. Its entire grade is $I+U\cdot C_X$, including every case tag and unused shape. Define
\[
 T_P^{\mathsf F}(X)=J_H(P)\odot\operatorname{Case}_U^{\mathsf F}(X),\qquad
 T_P^{\mathrm{old}}(Y)=P;\operatorname{Case}_U^{\mathrm{old}}(Y).
\]
The first uses all admitted child couplings and the second uses the original independent composition. Their full grade functor is the same $G_{C,U}$; they differ in their admitted law families, not in the whole case collection.

Selection of the distinguished member commutes with the typed case, and the distinguished mixed composite is the original dependent product. Therefore there is a natural full native isomorphism
\[
 \sigma_X:D_H T_P^{\mathsf F}(X)\xrightarrow{\sim}
                  T_P^{\mathrm{old}}D_H(X).             \tag{3}
\]
Its raw map is the dependent-tuple identification, removing only forced singleton $B$-identity witnesses; its coefficient map is the actual tensor/unit comparison. Its entire grade map is the corresponding identity on $G_{C,U}(C_X)$, and its body retains every $I$ versus $(u,d)$ alternative. Native parameter maps act on the same root/child factors on both routes, proving naturality. There is no claim that $T_P^{\mathsf F}J_H$ is isomorphic to $J_HT_P^{\mathrm{old}}$; its extra coupling laws generally prohibit that.

Let $(F_P,\iota)$ be the whole-case event-free family feedback algebra on the all-raw well-founded domain, and let $(R_P,a)$ be the old finite independent recursive process with its ordinary initial-algebra constructor. Write $\rho:D_HF_P\to R_P$ for (2) on loop inputs. Then
\[
 \rho\,D_H(\iota)\,\sigma_{F_P}^{-1}
       =a\,T_P^{\mathrm{old}}(\rho).                    \tag{4}
\]
Indeed each route grafts the same complete raw initializer tree, with the same independent product probability and actual tensor factors. On every formal whole-case shape it uses the same root constructor of $R_UC$. Thus (4) holds on all raw atoms, coefficients and unused grade maps, not only at selected states. It identifies $D_HF_P$ with the old initial algebra. This is a proved preservation result for this specific tree construction, not a general claim that right adjoints preserve initial algebras.

For any family algebra $\alpha:T_P^{\mathsf F}X\to X$, its distinguished interpretation is the old algebra
\[
 \alpha^{\mathrm{old}}=D_H(\alpha)\,\sigma_X^{-1}:
                   T_P^{\mathrm{old}}D_HX\to D_HX.
\]
The unique family fold $f:F_P\to X$ satisfies
\[
 D_H(f)\,\rho^{-1}:R_P\to D_HX
\]
equal to the unique old algebra fold into $\alpha^{\mathrm{old}}$. Applying $D_H$ to the family fold equation, using naturality of (3) and (4), proves that it is an old algebra morphism; old initiality supplies uniqueness. Equivalently the raw initializer-tree and full grade recursions on both routes are the same. This verifies fold compatibility, not merely preservation of a fixed-point carrier.

Finally the coreflection counit gives an actual native comparison
\[
 \kappa_P:J_HR_P\xrightarrow{J_H(\rho^{-1})}J_HD_HF_P
                         \xrightarrow{\varepsilon_{F_P}}F_P. \tag{5}
\]
It is identity on complete raw initializer trees, actual context factors, the entire $R_UC$ grade and bodies, and includes the distinguished independent law in the admitted family. It is natural for the fixed-$U$ source comparisons on the common all-raw guard domain. It need not be invertible: let $U=\{u_0,u_1\}$, let the root at $u_0$ deterministically emit two $u_1$ ports, and let the initializer at $u_1$ draw a fair raw bit and return one $B$ output with that bit. The only possible loop edge is $u_0\leadsto u_1$, so the guard holds. At $u_0$ the family feedback has the full four raw bit pairs and admits the product, diagonal and antidiagonal fair couplings; $J_HR_P$ admits only the product. All coefficient and full grade data can be retained identically. The family cardinality/identity-raw-law condition prevents a native inverse to (5).

Crucially, (5) is not automatically a morphism of family $T_P^{\mathsf F}$ algebras: $J_HR_P$ has not been supplied such an algebra. Its constructor domain can contain extra child couplings that no map to its singleton law family preserves. The algebra comparison established is instead (4) after applying $D_H$, and the compatible fold statement derived from it. This distinction prevents the directed source comparison from being mistaken for a strong family-feedback inclusion.
\end{document}
